\appendix

\subsection*{Rappels théoriques}

\paragraph{Moyenne superficielle :}
On considère une grandeur $g$ définie uniquement dans le volume fluide $V_l$.
On définit sa moyenne sur le volume total $V$ (fluide + solide) par :

\begin{equation}
\avg{g} = \frac{1}{V}\int_{V_l} g\, dV \label{eq:moy_sup}
\end{equation}
Cette moyenne, dite superficielle, s’exprime comme la moyenne dans le fluide multipliée 
par la porosité $\varepsilon_l = \frac{V_l}{V}$.\cite{le<3MultiscaleModelling2022}
Elle ne correspond donc pas à la moyenne de g dans le fluide seul.
La moyenne intrinsèque $\avgl{g}$ représente la vraie valeur moyenne de $g$ 
dans le fluide seul. La relation entre moyenne superficielle et moyenne 
intrinsèque est alors :

\begin{equation}
\avg{g} = \varepsilon_l \avgl{g} = \frac{V_l}{V}\avgl{g}  \label{eq:relation_moy}
\end{equation}

\paragraph{Théorème de la divergence de Gauss :}
Ce théorème relie une intégrale de volume à une intégrale de surface. 
Pour une quantité $g$ définie dans un volume $V$ délimité par une surface $A$ :

\begin{equation}
\int_V \nab g\, dV = \oint_A \mathbf{n}\, g\, dA
\label{eq:gauss}
\end{equation}

où $\mathbf{n}$ est le vecteur normal sortant à la surface $A$.

\paragraph{Théorème de la moyenne spatiale :}
Ce théorème découle du théorème de la divergence 
de Gauss appliqué au volume fluide.
Il relie la moyenne du gradient d'une quantité $g$ à son gradient 
de moyenne, avec un terme de correction surfacique sur l'interface 
$Asf$ entre le fluide et le solide :

\begin{equation}
\avg{\nab g} = \nab\avg{ g} + \frac{1}{V}\int_{Asf} \mathbf{n}\, g\, dA \label{eq:moy_spatiale}
\end{equation}

où $\mathbf{n}$ est le vecteur normal à l'interface $Asf$ 
pointant du fluide vers le solide.

\paragraph{Théorème de transport :}
Ce théorème permet d'échanger la moyenne et la dérivée temporelle, 
à condition que le volume $V_l$ ne dépende pas du temps :

\begin{equation}
\left\langle \frac{\partial g}{\partial t} \right\rangle = \frac{\partial}{\partial t}\avg{g} \label{eq:transport}
\end{equation}

\bigskip

% ============================================================
